\chapter{Mem\'oria virtual, demand paging e interrup\c{c}\~oes}
\label{cap:memoria-usuario}

Este cap\'itulo documenta os m\'odulos \code{gdt.rs}, \code{idt.rs},
\code{paging.rs} e \code{user.rs}, adicionados depois do console e do shell.
O objetivo \'e sair do modelo ``ring 0, interrup\c{c}\~oes desligadas e
mem\'oria est\'atica'' para um kernel capaz de:

\begin{itemize}
  \item receber exce\c{c}\~oes e IRQs por uma IDT;
  \item receber o timer do PIT em IRQ0;
  \item alocar frames f\'isicos de 4 KiB;
  \item materializar p\'aginas somente quando a CPU gera \code{\#PF};
  \item executar c\'odigo em ring 3 com uma barreira real entre usu\'ario e
        kernel.
\end{itemize}

\section{Ordem de inicializa\c{c}\~ao}

A sequ\^encia em \code{kernel\_main} passou a ser:

\begin{enumerate}
  \item \code{vga::init()} e \code{serial::init()};
  \item \code{paging::init(magic, info)} --- libera os frames dispon\'iveis
        descritos pelo Multiboot2;
  \item \code{gdt::init()} --- GDT definitiva, segmentos ring 3 e TSS;
  \item \code{keyboard::init()} --- teclado PS/2 por \emph{polling};
  \item \code{idt::init()} --- IDT, remapeamento do PIC e PIT;
  \item \code{sti} e entrada no shell.
\end{enumerate}

A IDT s\'o \'e instalada depois que a GDT tem a TSS, pois a porta de
\emph{double fault} usa uma pilha IST. O timer s\'o \'e desmascarado depois
que a IDT est\'a pronta.

\section{GDT, TSS e pilhas de privil\'egio}

A GDT do stub de boot cont\'em apenas descritores ring 0. Para entrar em
ring 3, \code{gdt.rs} monta uma tabela com sete entradas:

\begin{table}[H]
\centering
\caption{Descritores da GDT definitiva.}
\label{tab:gdt-usuario}
\begin{tabular}{@{}rll@{}}
\toprule
\textbf{Selector} & \textbf{Conte\'udo} & \textbf{DPL} \\
\midrule
\code{0x08} & c\'odigo 64 bits & 0 \\
\code{0x10} & dados & 0 \\
\code{0x18} & dados & 3 \\
\code{0x20} & c\'odigo 64 bits & 3 \\
\code{0x28} & TSS (duas entradas) & 0 \\
\bottomrule
\end{tabular}
\end{table}

A TSS \'e obrigat\'oria para a troca autom\'atica de pilha. Quando a CPU
passa de ring 3 para ring 0 por interrup\c{c}\~ao ou exce\c{c}\~ao, ela l\^e
\code{rsp0} da TSS e continua a empilhar nessa pilha. O kernel reserva:

\begin{itemize}
  \item \code{USER\_KERNEL\_STACK}: 16 KiB para interrup\c{c}\~oes vindas do
        usu\'ario;
  \item \code{DOUBLE\_FAULT\_STACK}: 8 KiB como \code{IST1}, para que um
        \code{\#DF} n\~ao tente usar uma pilha possivelmente corrompida.
\end{itemize}

O campo \code{iomap\_base} aponta para o fim da TSS, desativando o bitmap de
I/O. Assim, instru\c{c}\~oes \code{in}/\code{out} em ring 3 geram \code{\#GP}.

\section{IDT, stubs e normaliza\c{c}\~ao do erro}

A IDT tem 256 portas de 16 bytes. As 32 primeiras s\~ao exce\c{c}\~oes; as
16 seguintes s\~ao as IRQs do PIC remapeadas para \code{0x20..0x2F}; a porta
\code{0x80} \'e uma \emph{syscall} did\'atica com DPL = 3.

O problema cl\'assico \'e que a CPU empilha um c\'odigo de erro apenas em
algumas exce\c{c}\~oes (8, 10--14, 17, 21, 29 e 30). Se cada \emph{stub}
tratasse a pilha de forma diferente, o layout visto pelo Rust mudaria. Para
uniformizar, \code{idt.rs} gera stubs em assembly:

\begin{lstlisting}[caption={Stubs que normalizam o c\'odigo de erro.},label={lst:isrstubs}]
.macro isr_no_err n
.global isr_stub_\n
isr_stub_\n:
    pushq $0          # erro sintetico
    pushq $\n         # numero do vetor
    jmp isr_common
.endm

.macro isr_err n
.global isr_stub_\n
isr_stub_\n:
    pushq $\n         # erro real ja empilhado pela CPU
    jmp isr_common
.endm
\end{lstlisting}

O \emph{stub} comum salva os registradores de uso geral, preserva
\code{xmm0..xmm15} e \code{MXCSR}, alinha a pilha para chamar
\code{interrupt\_dispatch} e, no retorno, restaura tudo e executa
\code{iretq}. O frame visto pelo Rust \'e:

\begin{lstlisting}[caption={Layout de \code{InterruptFrame} ap\'os o stub comum.},label={lst:frame}]
#[repr(C)]
pub struct InterruptFrame {
    pub r15: u64, /* ... r8 */
    pub rdi: u64, /* ... rax */
    pub vector: u64,
    pub error: u64,
    pub rip: u64,
    pub cs: u64,
    pub rflags: u64,
    pub rsp: u64,   // so existe se veio de ring 3
    pub ss: u64,    // so existe se veio de ring 3
}
\end{lstlisting}

O vetor 14 chama o tratador de \emph{page fault}; o vetor 13 trata
\code{\#GP}; o vetor 0x80 encerra o programa de usu\'ario; o vetor 32
incrementa o contador de \emph{ticks} e envia EOI ao PIC.

\section{PIC e PIT}

O PIC 8259 \'e remapeado para n\~ao colidir com as exce\c{c}\~oes:

\begin{lstlisting}[caption={Remapeamento do PIC e inicializa\c{c}\~ao do PIT.},label={lst:picpit}]
outb(0x20, 0x11); outb(0xA0, 0x11);  // ICW1
outb(0x21, 0x20); outb(0xA1, 0x28);  // ICW2: vetores 0x20/0x28
outb(0x21, 0x04); outb(0xA1, 0x02);  // ICW3
outb(0x21, 0x01); outb(0xA1, 0x01);  // ICW4

let divisor = 11_931;                // 1,193182 MHz / 100 Hz
outb(0x43, 0x36);
outb(0x40, divisor as u8);
outb(0x40, (divisor >> 8) as u8);

outb(0x21, 0xFE);                    // so IRQ0
outb(0xA1, 0xFF);
\end{lstlisting}

A IRQ1 continua mascarada porque \code{keyboard.rs} l\^e o 8042 por
\emph{polling}. Habilit\'a-la exigiria um buffer de scancodes e coordenar o
acesso ao controlador.

\section{Alocador de frames f\'isicos}

\code{paging::init} percorre as tags Multiboot2. A tag 6 (mapa de mem\'oria)
fornece entradas de 24 bytes com base, tamanho e tipo. O alocador come\c{c}a
com todos os 262.144 frames do primeiro GiB marcados como usados e libera
somente as faixas do tipo 1 (RAM dispon\'ivel).

O bitmap tem 32 KiB:

\begin{lstlisting}[caption={Bitmap de frames e aloca\c{c}\~ao.},label={lst:bitmap}]
static mut FRAME_BITMAP: [u8; 32 * 1024] = [0xFF; 32 * 1024];

pub fn alloc_frame() -> Option<u64> {
    /* procura um bit zero, marca como usado e zera os 4096 bytes */
}
\end{lstlisting}

O primeiro MiB e o fim do kernel (s\'imbolo \code{\_\_kernel\_end} do
\emph{linker script}) ficam sempre reservados. A estrutura Multiboot2 tamb\'em
\'e reservada depois que as faixas dispon\'iveis s\~ao liberadas.

Como o stub mapeia o primeiro GiB por identidade, escrever no endere\c{c}o
f\'isico do frame \'e equivalente a escrever no endere\c{c}o virtual baixo.
Isso permite zerar frames sem precisar de um mapa direto adicional.

\section{Demand paging e prote\c{c}\~ao U/S}

Cada entrada de tabela tem bits de presen\c{c}a (P), escrita (W) e
usu\'ario (U/S). O x86\_64 exige U/S = 1 em \emph{todos} os n\'iveis
(PML4, PDPT, PD e PT) para permitir um acesso em ring 3. O kernel usa isso
para criar duas regi\~oes:

\begin{table}[H]
\centering
\caption{Regi\~oes de demand paging.}
\label{tab:demand}
\begin{tabularx}{\linewidth}{@{}lrrX@{}}
\toprule
\textbf{Regi\~ao} & \textbf{Base} & \textbf{Tamanho} & \textbf{U/S das folhas} \\
\midrule
Usu\'ario & \code{0x0000008000000000} & 2 MiB & 1 \\
Kernel   & \code{0xFFFFFFFFC0000000} & 2 MiB & 0 \\
\bottomrule
\end{tabularx}
\end{table}

Nenhuma folha dessas regi\~oes existe no boot. Quando o \code{\#PF} ocorre,
o tratador l\^e \code{CR2} e o c\'odigo de erro e decide:

\begin{enumerate}
  \item se o endere\c{c}o pertence \`a regi\~ao de usu\'ario e a p\'agina
        n\~ao est\'a presente, aloca um frame e mapeia com
        \code{P|W|U/S}; as tabelas intermedi\'arias tamb\'em recebem U/S = 1;
  \item se pertence \`a regi\~ao do kernel e o bit U/S do erro \'e zero,
        mapeia com \code{P|W} e U/S = 0;
  \item se um usu\'ario tentar a regi\~ao do kernel, o erro tem U/S = 1 e o
        tratador recusa o mapeamento; o programa de usu\'ario \'e encerrado;
  \item se o endere\c{c}o est\'a fora das duas regi\~oes, a falha \'e fatal
        para o kernel ou encerra o usu\'ario.
\end{enumerate}

\begin{nota}[Por que o kernel n\~ao pode ``abrir'' uma tabela do supervisor]
Mapear uma p\'agina de usu\'ario n\~ao consiste apenas em ligar U/S na folha.
Se qualquer n\'ivel intermedi\'ario for supervisor, o hardware nega o acesso.
Al\'em disso, transformar uma tabela do kernel em tabela de usu\'ario poderia
expor outros mapeamentos. Por isso \code{next\_table} recusa uma tabela
supervisor no caminho de uma p\'agina de usu\'ario.
\end{nota}

\section{Modo usu\'ario e retorno ao kernel}

\code{user.rs} monta um pequeno programa de m\'aquina em um buffer e o copia
para a p\'agina \code{USER\_CODE}, que \'e residente. A pilha e a p\'agina de
dados ficam sob demanda:

\begin{enumerate}
  \item o programa faz \code{push}/\code{pop}, gerando \code{\#PF} na pilha;
  \item escreve em \code{USER\_DATA}, gerando outro \code{\#PF};
  \item executa \code{int \$0x80} com \code{rax = 0}.
\end{enumerate}

O comando \code{userprot} usa outro programa: ele tenta ler
\code{0xFFFFFFFF80000000} (in\'icio do kernel). A p\'agina existe, mas tem
U/S = 0; o erro de \emph{page fault} tem P = 1 e U/S = 1. O tratador
reconhece a viola\c{c}\~ao, imprime os dados e n\~ao devolve o controle ao
programa.

A entrada em ring 3 \'e feita por \code{iretq}, que carrega CS/SS de DPL 3,
RIP, RSP e RFLAGS. A volta n\~ao usa o caminho normal do \code{iretq}:

\begin{itemize}
  \item \code{user\_enter} salva os registradores callee-saved e o RSP do
        kernel em \code{user\_saved\_rsp};
  \item uma exce\c{c}\~ao ou \code{int \$0x80} cai no \emph{stub} comum em
        ring 0;
  \item \code{user::abort} grava o status e chama \code{user\_resume}, que
        restaura o RSP salvo, recarrega os segmentos do kernel e retorna para
        a chamada de \code{user\_enter}.
\end{itemize}

Esse mecanismo \'e um \emph{longjmp} controlado. Ele evita ter de reconstruir
o frame de interrup\c{c}\~ao para voltar a ring 0 e mant\'em o shell simples.

\section{Comandos e evid\^encias}

\begin{table}[H]
\centering
\caption{Comandos que exercitam a nova infraestrutura.}
\label{tab:comandos-memoria}
\begin{tabularx}{\linewidth}{@{}lX@{}}
\toprule
\textbf{Comando} & \textbf{Efeito} \\
\midrule
\code{demand} & escreve na regi\~ao de demand paging do kernel; o primeiro
acesso gera \code{\#PF} \\
\code{usertest} & entra em ring 3, gera \code{\#PF} na pilha e nos dados e
sai por \code{int \$0x80} \\
\code{userprot} & entra em ring 3 e tenta ler o kernel; a prote\c{c}\~ao U/S
encerra o programa \\
\code{mem} & mostra frames alocados/livres, \emph{page faults}, \emph{ticks}
e os bits U/S de uma p\'agina de usu\'ario e do kernel \\
\bottomrule
\end{tabularx}
\end{table}

Uma execu\c{c}\~ao t\'ipica em QEMU:

\begin{lstlisting}[style=shell,caption={Sa\'ida dos comandos de mem\'oria e ring 3.},label={lst:testdemand}]
rinux> demand
[PF] demanda atendida: 0xffffffffc0000000 (kernel, erro=0x2)
[DEMAND] 0xffffffffc0000000 = 0xdeadbeefcafe1234 (page faults: 1)
rinux> usertest
[USER] entrando em ring 3 (demand paging)...
[PF] demanda atendida: 0x0000008000007ff8 (usuario, erro=0x6)
[PF] demanda atendida: 0x0000008000002000 (usuario, erro=0x6)
[USER] voltou ao kernel; page faults nesta execucao: 2
rinux> userprot
[USER] entrando em ring 3 (teste de protecao)...
[USER] page fault em 0xffffffff80000000: endereco fora das regioes de demand paging (erro=0x5, rip=0x000000800000100a)
[USER] programa encerrado com status 1
rinux> mem
Frames fisicos (4 KiB):
  alocados : 7
  livres   : 32217
  page faults atendidos: 3
  ticks do timer       : 372
  USER_CODE  flags=0x203067 (U/S=1)
  KERNEL_VMA flags=0xe3 (U/S=0)
\end{lstlisting}

O erro \code{0x6} (P = 0, W = 1, U/S = 1) confirma uma escrita de usu\'ario em
p\'agina ausente; o erro \code{0x5} (P = 1, U/S = 1, sem W) confirma uma
leitura de usu\'ario em p\'agina presente e supervisor.

\section{Limita\c{c}\~oes desta implementa\c{c}\~ao}

\begin{itemize}
  \item n\~ao h\'a espa\c{c}os de endere\c{c}o separados por processo: o
        programa de usu\'ario usa a mesma PML4 do kernel;
  \item n\~ao h\'a escalonador nem troca de contexto; o timer apenas conta
        \emph{ticks};
  \item a \'unica \emph{syscall} \'e \code{exit} (\code{rax = 0}); n\~ao
        h\'a valida\c{c}\~ao de ponteiros de usu\'ario;
  \item SMEP/SMAP e NX ainda n\~ao est\~ao habilitados; a separa\c{c}\~ao
        depende apenas de U/S;
  \item o teclado continua em \emph{polling} e IRQ1 permanece mascarada;
  \item as regi\~oes de demand paging t\^em tamanho fixo de 2 MiB e ficam
        fora do primeiro GiB mapeado por identidade.
\end{itemize}
